\section{Object Detection}
In order to localize the ski jumper within each video frame, we employed an object detection technique based on convolutional neural networks. Specifically, we selected the YOLOv8 architecture for this task. 

\subsection{Fine-tuning}
To acquire the optimal model for ski jumper detection, we utilized fine-tuning to adapt a pre-trained YOLOv8 model. The goal was to recognize the jumper's bounding box with the highest possible accuracy.

This task required a well-defined and precise dataset. To construct this, we extracted frames from experimental videos at random time intervals ranging from 500 ms to 2500 ms. This sampling strategy ensured diverse temporal coverage of the entire jump sequence. As a result of this extraction process, we obtained a base dataset of 991 frames.

Following extraction, we proceeded with data annotation using the \href{https://app.roboflow.com/skijumpinglabeling/skijumping/2}{Roboflow} platform. Bounding boxes were drawn tightly around the jumper to minimize background noise and exclude unnecessary margins. Additionally, we included 16 unannotated frames (background only) to improve the model's robustness against false positives.

To prevent overfitting and improve the model's generalization capabilities regarding angle and lighting conditions, we applied data augmentation. Specifically, we utilized rotation (between -15\textdegree{} and +15\textdegree{}) and brightness adjustment (between -15\% and +15\%). 

For training and validation purposes, the augmented dataset was split as follows:
\begin{itemize}
    \item \textbf{Train set:} 2294 images (91\%)
    \item \textbf{Valid set:} 159 images (6\%)
    \item \textbf{Test set:} 67 images (3\%)
\end{itemize}

With the dataset prepared, we proceeded with fine-tuning using the following configuration:
\begin{enumerate}
    \item \textbf{Model Selection:} Initialization of the pre-trained YOLOv8 Nano (\texttt{yolov8n.pt}) as the starting checkpoint.
    \item \textbf{Training Execution:} The training was executed via the \texttt{model.train()} function with 100 epochs, an image size of 640 pixels, using an NVIDIA T4 GPU (Device 0).
\end{enumerate}

The process yielded a highly effective model, confirmed by the final performance metrics:
\begin{itemize}
    \item \textbf{mAP50 (0.992):} Indicates that the model almost always detects the jumper accurately.
    \item \textbf{mAP50-95 (0.817):} Suggests that the predicted bounding boxes align precisely with the jumper and their equipment.
    \item \textbf{Precision (0.983) and Recall (0.975):} Confirms that the model produces very few false positives and successfully identifies the jumper in nearly every frame.
\end{itemize}

\subsection{Detection Pipeline}
Leveraging the fine-tuned model, we constructed a pipeline for continuous detection and tracking across video frames.

The pipeline processes the video frame-by-frame. First, the YOLO model localizes the jumper. In cases where the model detects multiple targets, the system selects the bounding box with the highest confidence score. To maintain continuity, we utilized YOLO's object tracking capabilities, assigning a unique ID to the detected subject to distinguish the continuous trajectory from momentary detection artifacts.

To determine the jumper's precise Center of Mass (CoM), we integrated the Segment Anything Model (SAM). SAM is an advanced segmentation model capable of separating the jumper from the background with high pixel-level precision. Since SAM is computationally intensive, it is applied only within the bounding box coordinates provided by YOLO (the Region of Interest). Using the precise segmentation mask generated by SAM, the model calculates the exact geometric center of the jumper.

However, the raw CoM coordinates calculated frame-by-frame often exhibit high-frequency noise or "jitter" due to minor segmentation fluctuations. To address this, we implemented a Kalman filter to estimate and smooth the jumper's trajectory.  The Kalman filter predicts the future position of the point based on previous states and corrects it using the current measurement. Consequently, the final tracking point moves smoothly, as the filter effectively mitigates noise by statistically weighing the predicted state against the measured position.

\section{Trajectory estimation}

\section{Homography from camera plane to the scene}